
\documentclass[11pt]{article}
\usepackage[utf8]{inputenc}
\usepackage[T1]{fontenc}
\usepackage[most]{tcolorbox}
\usepackage{amsmath, amssymb}
\usepackage{bm}
\usepackage{graphicx}
\usepackage{listings}
\usepackage{enumitem}
\usepackage{lmodern}
\usepackage{xcolor}
\usepackage{float}
\usepackage{color}
\usepackage{fancyvrb}
\usepackage{geometry}
\geometry{margin=1in}

\definecolor{codegray}{rgb}{0.5,0.5,0.5}
\definecolor{codepurple}{rgb}{0.58,0,0.82}
\definecolor{backcolour}{rgb}{0.95,0.95,0.92}

\lstdefinestyle{mystyle}{
    backgroundcolor=\color{backcolour},
    commentstyle=\color{codegray},
    keywordstyle=\color{blue},
    numberstyle=\tiny\color{gray},
    stringstyle=\color{codepurple},
    basicstyle=\ttfamily\footnotesize,
    breakatwhitespace=false,
    breaklines=true,
    captionpos=b,
    keepspaces=true,
    numbers=left,
    numbersep=5pt,
    showspaces=false,
    showstringspaces=false,
    showtabs=false,
    tabsize=2
}

\lstset{style=mystyle}
\title{\textbf{Generative Modeling with Optimal Transport Maps}\\[0.3em]\large Paper review: Rout, Burnaev \& Korotin, ICLR 2022}
\author{Pierre Chambet}
\date{May 2025}

\begin{document}
\maketitle

\section*{Key Formulas: AE-OT vs OT-map (ICLR 2022)}

This summary compares the core mathematical formulations of AE-OT (Autoencoder + Optimal Transport) and the OT-map method from ``Generative Modeling with Optimal Transport Maps'' (ICLR 2022).

\subsection*{\textbf{AE-OT: Autoencoder + Optimal Transport}}

\textbf{Phase 1 – Autoencoder training (supervised reconstruction):}
\[
\min_{E, D} \; \mathbb{E}_{y \sim \nu} \left[ \| D(E(y)) - y \|^2 \right]
\]
\begin{itemize}
    \item $y \in \mathbb{R}^D$: real image
    \item $z = E(y) \in \mathbb{R}^H$: latent representation
\end{itemize}

\textbf{Phase 2 – Discrete OT between noise and latent codes:} \\
Let $x_i \sim \mu$ (noise), and $z_j = E(y_j) \sim E_\# \nu$.

\textbf{Option 1: Monge OT (with permutation):}
\[
\min_T \sum_{i=1}^n \| T(x_i) - z_{\pi(i)} \|^2
\]

\textbf{Option 2: Entropic Kantorovich OT with barycentric projection:}
\[
\min_P \sum_{i,j} P_{ij} \| x_i - z_j \|^2 + \varepsilon \sum_{i,j} P_{ij} \log P_{ij}
\]
\[
T(x_i) = \sum_j P_{ij} z_j \quad \text{(barycenter)}
\]

\textbf{Final generation:}
\[
\hat{y} = D(T(x))
\]

\subsection*{\textbf{OT-map (ICLR 2022)}}

\textbf{Setting:} Learn transport map $G(x) = T(Q(x))$ directly in the image space $\mathbb{R}^D$, without any encoder.

\textbf{Min--max formulation (dual of OT):}
\[
\min_{\psi} \max_{G} \left[ \mathbb{E}_{x \sim \mu} \left( \langle Q(x), G(x) \rangle - \psi(G(x)) \right) + \mathbb{E}_{y \sim \nu} \psi(y) \right]
\]
\begin{itemize}
    \item $Q(x) \in \mathbb{R}^D$: embedded noise (e.g., via upsampling)
    \item $G(x) \in \mathbb{R}^D$: image generated from noise
    \item $\psi(y)$: Kantorovich potential (convex function)
\end{itemize}

\textbf{Optimal transport geometry:}
\[
G(x) \approx \arg\max_y \left( \langle Q(x), y \rangle - \psi(y) \right)
\quad \Rightarrow \quad G(x) \in \partial \psi^*(Q(x))
\]

\textbf{No decoder needed — generation is direct:}
\[
\hat{y} = G(x)
\]

\subsection*{Comparison}

\begin{center}
\begin{tabular}{|l|l|l|}
\hline
\textbf{Feature} & \textbf{AE-OT} & \textbf{OT-map (this paper)} \\
\hline
Space of transport & Latent space $\mathbb{R}^H$ & Ambient space $\mathbb{R}^D$ \\
Needs autoencoder? & Yes & No \\
Uses encoder/decoder? & Yes: $E$, $D$ & No \\
Transport computed via & Discrete OT (Monge or Sinkhorn) & Min--max dual OT formulation \\
Decoding step & $\hat{y} = D(T(x))$ & $\hat{y} = G(x)$ (direct output) \\
Type of supervision & Point set alignment & Distribution-level alignment \\
\hline
\end{tabular}
\end{center}


\section*{Derivation of the OT-map Objective}

\subsection*{1. Monge’s Problem (Primal Formulation)}

Let \(\mu\) and \(\nu\) be two probability measures defined on metric spaces \(\mathcal{X}, \mathcal{Y} \subset \mathbb{R}^d\). The original problem posed by Monge (1781) is:

\[
\inf_{T_\#\mu = \nu} \int_{\mathcal{X}} c(x, T(x)) \, d\mu(x)
\]

where \(T : \mathcal{X} \rightarrow \mathcal{Y}\) is a measurable map that pushes forward \(\mu\) to \(\nu\) (i.e., \(T_\#\mu = \nu\)), and \(c(x, y)\) is a ground cost function.

\textbf{Remark:} Monge’s problem may be ill-posed; it does not allow mass splitting and may lack a solution if, e.g., \(\mu\) is a Dirac delta and \(\nu\) is not.

\subsection*{2. Kantorovich’s Relaxation}

Kantorovich (1942) relaxed the problem by allowing probabilistic couplings \(\pi\) between \(\mu\) and \(\nu\). The new problem becomes:

\[
\inf_{\pi \in \Pi(\mu, \nu)} \int_{\mathcal{X} \times \mathcal{Y}} c(x, y) \, d\pi(x, y)
\]

where \(\Pi(\mu, \nu)\) is the set of all joint distributions with marginals \(\mu\) and \(\nu\).

\subsection*{3. Dual Formulation and $c$-Transform}

The Kantorovich dual problem is:

\[
\sup_{u, v} \left\{ \int_{\mathcal{X}} u(x) \, d\mu(x) + \int_{\mathcal{Y}} v(y) \, d\nu(y) \,:\, u(x) + v(y) \leq c(x, y) \right\}
\]

This is a linear program over functions \(u, v\). To work more effectively with this constraint, we define the \textbf{\(c\)-transform}:

\subsection*{Definition: \(c\)-Transform}

Given \(v : \mathcal{Y} \to \mathbb{R}\), the \(c\)-transform of \(v\) is:

\[
v^c(x) := \inf_{y \in \mathcal{Y}} \left\{ c(x, y) - v(y) \right\}
\]

Analogously, the \(c\)-transform of \(u\) is:

\[
u^c(y) := \inf_{x \in \mathcal{X}} \left\{ c(x, y) - u(x) \right\}
\]

Using the \(c\)-transform, we can rewrite the dual in terms of a single potential:

\[
\sup_v \left\{ \int_{\mathcal{X}} v^c(x) \, d\mu(x) + \int_{\mathcal{Y}} v(y) \, d\nu(y) \right\}
\]

This is equivalent to:

\[
\sup_u \left\{ \int_{\mathcal{X}} u(x) \, d\mu(x) + \int_{\mathcal{Y}} u^c(y) \, d\nu(y) \right\}
\]

\subsection*{4. Specialization to Quadratic Cost and Convex Potentials}

Let us now consider the quadratic cost (the OT value is then $\frac{1}{2} W_2^2$):

\[
c(x, y) = \frac{1}{2} \|x - y\|^2
\]

\textbf{Why convex functions?} Brenier's theorem (1991) states: if \(\mu\) is absolutely continuous and \(c(x, y) = \frac{1}{2} \|x - y\|^2\), then the optimal transport map \(T^*\) is the gradient of a convex function:

\[
T^*(x) = \nabla \varphi(x)
\]

This justifies parameterizing the OT map as the gradient of a convex potential.

\subsection*{5. Reformulating the Dual using Convex Potentials}

Let us define:

\[
\psi(y) := \frac{1}{2} \|y\|^2 - v(y)
\]

Then we obtain:

\[
v(y) = \frac{1}{2} \|y\|^2 - \psi(y)
\]

and so:


\[
\begin{aligned}
v^c(x) &= \inf_y \left\{ \frac{1}{2} \|x - y\|^2 - v(y) \right\} \\
&= \inf_y \left\{ \frac{1}{2} \|x - y\|^2 - \left( \frac{1}{2} \|y\|^2 - \psi(y) \right) \right\}\\
&= \inf_y \left\{ \frac{1}{2} \|x\|^2 - \langle x, y \rangle + \psi(y) \right\}\\
&= \frac{1}{2} \|x\|^2 - \sup_y \left\{ \langle x, y \rangle - \psi(y) \right\}\\
&= \frac{1}{2} \|x\|^2 - \psi^*(x)
\end{aligned}
\]

Thus, the dual becomes:

\[
\frac{1}{2} W_2^2(\mu, \nu) = \int \frac{1}{2} \|x\|^2 \, d\mu(x) + \int \frac{1}{2} \|y\|^2 \, d\nu(y) - \inf_{\psi \text{ convex}} \left\{ \int \psi^*(x) \, d\mu(x) + \int \psi(y) \, d\nu(y) \right\}
\]

\subsection*{6. Final Form: Saddle Point Optimization}

Using the Fenchel–Legendre duality:

\[
\psi^*(x) = \sup_{y \in \mathbb{R}^d} \left\{ \langle x, y \rangle - \psi(y) \right\}
\]

So:

\[
\int \psi^*(x) \, d\mu(x) = \sup_{T} \int \langle x, T(x) \rangle - \psi(T(x)) \, d\mu(x)
\]

The final expression is:

\[
\frac{1}{2} W_2^2(\mu, \nu) = \int \frac{1}{2} \|x\|^2 \, d\mu(x) + \int \frac{1}{2} \|y\|^2 \, d\nu(y)
- \inf_{\psi \text{ convex}} \sup_T \left\{ \int \langle x, T(x) \rangle - \psi(T(x)) \, d\mu(x) + \int \psi(y) \, d\nu(y) \right\}
\]

\subsection*{7. Optimization Objective for Generative Modeling}

Letting the constant terms be absorbed, we minimize the functional:

\[
\min_{\psi} \max_{T} \left\{ \int \langle x, T(x) \rangle - \psi(T(x)) \, d\mu(x) + \int \psi(y) \, d\nu(y) \right\}
\]

This is the saddle point problem solved in the ICLR 2022 paper, where both \(T\) and \(\psi\) are parameterized by neural networks.

\end{document}
